\paragrafos{Δυνάμεις ρητών: οι πολλαπλασιασμοί που επαναλαμβάνονται}{Α.3.6, «Δυνάμεις ρητών με εκθέτη μη μηδενικό φυσικό»}

\begin{anagke}
Στην αυλή του σχολείου ο Νίκος αφήνει μια μπάλα του μπάσκετ από ύψος $2$ μέτρων. Κάθε φορά που αναπηδά, φτάνει στο μισό του ύψους της προηγούμενης αναπήδησης. «Σε τι ύψος θα φτάσει στην τρίτη αναπήδηση;» ρωτάει η Πολυξένη. Ο Ηλίας γράφει: μετά την πρώτη $2\cdot\frac{1}{2}$, μετά τη δεύτερη $2\cdot\frac{1}{2}\cdot\frac{1}{2}$, μετά την τρίτη $2\cdot\frac{1}{2}\cdot\frac{1}{2}\cdot\frac{1}{2}$. «Αυτό γράφεται πιο σύντομα», λέει η Ρήνη: «το $\frac{1}{2}$ τρεις φορές». Στους φυσικούς το ξέρουν ήδη: $2\cdot2\cdot2=2^3$. Μπορεί να γίνει και με κλάσματα;

Η κυρία Αλεξάνδρα βάζει στον πίνακα δύο παραστάσεις: $\left(-\frac{1}{2}\right)^3$ και $-\frac{1}{2^3}$. «Είναι ίδιες;» ρωτά. Ο Στέλιος λέει «ναι, όλα είναι μείον ένα όγδοο» και η Φανή γράφει στο τετράδιό της: «ίδια, αλλά γιατί;». Ο Βασίλης προτείνει $(-3)^2=-9$ και η Ζωή τον ρωτάει «πώς βγαίνει αρνητικό το γινόμενο δύο αρνητικών;»

Σε αυτή την παράγραφο θα επεκτείνουμε τις δυνάμεις στους ρητούς: θα μάθουμε να προσδιορίζουμε το πρόσημο μιας δύναμης και να υπολογίζουμε παραστάσεις με δυνάμεις.
\end{anagke}

\begin{lexi}{δύναμη, βάση, εκθέτης}{δύναμη, βάση, εκθέτης}
Η \lexiword{δύναμη} είναι η «ύψωση», το γινόμενο ίσων παραγόντων. Στα αρχαία ελληνικά «δύναμις» ήταν και το τετράγωνο. Η \lexiword{βάση} είναι ο αριθμός που πολλαπλασιάζεται με τον εαυτό του και ο \lexiword{εκθέτης} (από το «εκ-τίθημι», «βάζω έξω») ο αριθμός που γράφουμε πάνω δεξιά και δείχνει πόσους παράγοντες έχουμε.
\end{lexi}

\begin{metafrasi}
Το $\frac{2}{3}$ τρεις φορές. & $\left(\dfrac{2}{3}\right)^{3}=\dfrac{2}{3}\cdot\dfrac{2}{3}\cdot\dfrac{2}{3}=\dfrac{8}{27}$\\
Δύο αρνητικοί παράγοντες. & $(-1{,}5)^2=+2{,}25$\\
Τρεις αρνητικοί παράγοντες. & $(-0{,}1)^3=-0{,}001$\\
Η δύναμη με εκθέτη $1$. & $\alpha^1=\alpha$\\
Το μείον μόνο, χωρίς παρένθεση. & $-2^4=-16$\\
\end{metafrasi}

\tmima{Η θεωρία}

\begin{orismos}{δύναμη ρητού}
Η \textbf{δύναμη} του $\alpha$ στην $\nu$ ($\nu$ θετικός φυσικός) είναι το γινόμενο $\nu$ παραγόντων ίσων με $\alpha$:
$$\alpha^{\nu}=\underbrace{\alpha\cdot\alpha\cdot\ldots\cdot\alpha}_{\nu\ \text{παράγοντες}}.$$
Ο $\alpha$ λέγεται \textbf{βάση} (είναι ρητός) και ο $\nu$ \textbf{εκθέτης}. Ισχύει $\alpha^1=\alpha$.
\end{orismos}

Για παράδειγμα $\left(\frac{2}{3}\right)^3=\frac{2}{3}\cdot\frac{2}{3}\cdot\frac{2}{3}=\frac{2^3}{3^3}=\frac{8}{27}$. Γενικά, αφού το γινόμενο κλασμάτων έχει αριθμητή το γινόμενο των αριθμητών και παρονομαστή το γινόμενο των παρονομαστών, ισχύει $\left(\dfrac{\alpha}{\beta}\right)^{\nu}=\dfrac{\alpha^{\nu}}{\beta^{\nu}}$.

\begin{kanonas}{το πρόσημο μιας δύναμης}
\begin{itemize}[leftmargin=1.2em,itemsep=0.3ex,topsep=2pt]
\item Δύναμη με βάση \textbf{θετικό} αριθμό είναι πάντα \textbf{θετική}: αν $\alpha>0$, τότε $\alpha^{\nu}>0$.
\item Δύναμη με βάση \textbf{αρνητικό} αριθμό και εκθέτη \textbf{άρτιο} είναι \textbf{θετική}: αν $\alpha<0$ και $\nu$ άρτιος, τότε $\alpha^{\nu}>0$. Π.χ. $(-1{,}5)^2=(-1{,}5)\cdot(-1{,}5)=+2{,}25$.
\item Δύναμη με βάση \textbf{αρνητικό} αριθμό και εκθέτη \textbf{περιττό} είναι \textbf{αρνητική}: αν $\alpha<0$ και $\nu$ περιττός, τότε $\alpha^{\nu}<0$. Π.χ. $(-0{,}1)^3=-0{,}001$.
\end{itemize}
\end{kanonas}

Ο λόγος είναι ο κανόνας των προσήμων του γινομένου: κάθε ζευγάρι αρνητικών παραγόντων δίνει θετικό, άρα το πρόσημο εξαρτάται από το αν ο αριθμός των αρνητικών παραγόντων είναι άρτιος ή περιττός. Ειδικά $(-1)^{\nu}=1$ αν ο $\nu$ είναι άρτιος και $-1$ αν είναι περιττός. Επίσης $1^{\nu}=1$ και $0^{\nu}=0$.

\tmimasmall{Προτεραιότητα των πράξεων}

Στις παραστάσεις με πράξεις εκτελούμε με τη σειρά: πρώτα τις πράξεις μέσα σε \textbf{παρενθέσεις}, μετά τις \textbf{δυνάμεις}, μετά τους \textbf{πολλαπλασιασμούς και τις διαιρέσεις} (από αριστερά προς τα δεξιά) και τέλος τις \textbf{προσθέσεις και αφαιρέσεις}.

\begin{prosoxi}
\begin{itemize}[leftmargin=1.2em,itemsep=0.3ex,topsep=2pt]
\item Το $\alpha^{3}$ δεν είναι $3\cdot\alpha$, αλλά $\alpha\cdot\alpha\cdot\alpha$. Για παράδειγμα $0{,}2^4\neq4\cdot0{,}2$: $0{,}2^4=0{,}0016$.
\item Το $(-3)^2$ και το $-3^2$ δεν είναι ίδια. Το πρώτο είναι $(-3)(-3)=9$ και το δεύτερο είναι ο αντίθετος του $3^2$, δηλαδή $-9$. Η βάση είναι ό,τι βρίσκεται μέσα στην παρένθεση.
\item Στο $\left(\frac{2}{3}\right)^3$ ο εκθέτης δουλεύει και στον αριθμητή και στον παρονομαστή: $\frac{2^3}{3^3}=\frac{8}{27}$ και όχι $\frac{2^3}{3}$.
\item Ο εκθέτης δείχνει πλήθος παραγόντων, όχι πρόσθεση: $2^3=2\cdot2\cdot2=8$, όχι $2+2+2=6$.
\item Η δύναμη με άρτιο εκθέτη δεν είναι ποτέ αρνητική, ακόμη και με αρνητική βάση.
\end{itemize}
\end{prosoxi}

\tmima{Λυμένα παραδείγματα}

\paradeigma{Η μπάλα του Νίκου}
\tixero Η μπάλα αφήνεται από ύψος $2$ μ. και κάθε αναπήδηση φτάνει στο μισό της προηγούμενης.\par
\tithelo Σε τι ύψος φτάνει στην τρίτη αναπήδηση;\par
\possindeo Κάθε αναπήδηση πολλαπλασιάζει το ύψος επί $\frac{1}{2}$. Μετά από τρεις: $2\cdot\left(\frac{1}{2}\right)^3=2\cdot\frac{1}{8}=\frac{2}{8}=\frac{1}{4}$.\par
\apantisi $\frac{1}{4}$ μ., δηλαδή $25$ εκατοστά.\par
\elegxos Βήμα βήμα: $2\to1\to\frac{1}{2}\to\frac{1}{4}$.

\paradeigma{Δυνάμεις με πρόσημο}
\tixero Οι δυνάμεις $\left(\frac{2}{3}\right)^3$, $(-1{,}5)^2$, $(-0{,}1)^3$.\par
\tithelo Να τις υπολογίσω.\par
\possindeo Πρώτα το πρόσημο: $\frac{2}{3}$ θετική βάση, άρα θετικό. Το $(-1{,}5)^2$ έχει αρνητική βάση και άρτιο εκθέτη, άρα θετικό. Το $(-0{,}1)^3$ έχει αρνητική βάση και περιττό εκθέτη, άρα αρνητικό. Μετά οι απόλυτες τιμές: $\frac{8}{27}$, $1{,}5\cdot1{,}5=2{,}25$, $0{,}1\cdot0{,}1\cdot0{,}1=0{,}001$.\par
\apantisi $\dfrac{8}{27}$, $\ 2{,}25$, $\ -0{,}001$.\par
\elegxos Το $0{,}001$ έχει τρία δεκαδικά ψηφία: $1+1+1$.

\paradeigma{Η παρένθεση μετράει}
\tixero Οι $(-2)^4$, $-2^4$, $\left(-\frac{1}{2}\right)^3$ και $-\frac{1}{2^3}$.\par
\tithelo Να τις υπολογίσω και να τις συγκρίνω.\par
\possindeo Το $(-2)^4$ έχει αρνητική βάση και άρτιο εκθέτη: $(-2)\cdot(-2)\cdot(-2)\cdot(-2)=16$. Το $-2^4$ είναι ο αντίθετος του $2^4$: $-16$. Το $\left(-\frac{1}{2}\right)^3=-\frac{1}{8}$. Το $-\frac{1}{2^3}=-\frac{1}{8}$.\par
\apantisi $16$, $-16$, $-\dfrac{1}{8}$, $-\dfrac{1}{8}$. Οι δύο τελευταίες είναι ίσες.\par
\elegxos Το $(-2)^4$ είναι θετικό, ενώ το $-2^4$ αρνητικό: διαφέρουν.

\paradeigma{Η ειδική περίπτωση του $-1$}
\tixero Οι δυνάμεις $(-1)^{100}$, $(-1)^{101}$, $1^{50}$, $0^3$.\par
\tithelo Να τις υπολογίσω χωρίς να γράψω όλους τους παράγοντες.\par
\possindeo Το $(-1)^{\nu}$ είναι $1$ για άρτιο και $-1$ για περιττό $\nu$. Το $1^{\nu}=1$ και $0^{\nu}=0$.\par
\apantisi $1$, $-1$, $1$, $0$.\par
\elegxos Σε ζεύγη: $(-1)^{100}=\big((-1)(-1)\big)^{50}=1^{50}=1$.

\paradeigma{Παράσταση με προτεραιότητα}
\tixero Η παράσταση $3\cdot\left(\frac{1}{2}-\frac{5}{6}\right)+\left(\frac{1}{3}-\frac{2}{3}\right)^2$.\par
\tithelo Να βρω την τιμή της.\par
\possindeo Πρώτα οι παρενθέσεις: $\frac{1}{2}-\frac{5}{6}=\frac{3}{6}-\frac{5}{6}=-\frac{1}{3}$ και $\frac{1}{3}-\frac{2}{3}=-\frac{1}{3}$. Μετά η δύναμη: $\left(-\frac{1}{3}\right)^2=\frac{1}{9}$. Μετά ο πολλαπλασιασμός: $3\cdot\left(-\frac{1}{3}\right)=-1$. Τέλος η πρόσθεση: $-1+\frac{1}{9}=-\frac{9}{9}+\frac{1}{9}=-\frac{8}{9}$.\par
\apantisi $-\dfrac{8}{9}$.\par
\elegxos Με δεκαδικούς: $3\cdot(-0{,}333\ldots)+0{,}111\ldots=-0{,}888\ldots$

\tmima{Ασκήσεις}

\zoni{Το κατάλαβα}{ορισμός, πρόσημο, υπολογισμός}

\ex Σωστό ή λάθος; α) $\left(\frac{5}{2}\right)^2=\frac{5^2}{2^2}$. \quad β) $\left(\frac{2}{3}\right)^3=\frac{2^3}{3}$. \quad γ) $0{,}2^4=4\cdot0{,}2$. \quad δ) $\left(-\frac{3}{10}\right)^2=0{,}3^2$. \quad ε) $\left(-\frac{1}{3}\right)^3=-\left(\frac{1}{3}\right)^3$.

\ex Υπολόγισε: α) $\left(\frac{1}{2}\right)^2$ \quad β) $\left(-\frac{1}{3}\right)^3$ \quad γ) $\left(\frac{2}{5}\right)^2$ \quad δ) $(-0{,}2)^2$ \quad ε) $0{,}3^3$ \quad στ) $(-1)^7$.

\ex Βρες το πρόσημο χωρίς να υπολογίσεις: α) $\left(-\frac{3}{4}\right)^5$ \quad β) $(-0{,}2)^4$ \quad γ) $\left(-\frac{1}{2}\right)^{10}$ \quad δ) $-\left(\frac{3}{5}\right)^3$ \quad ε) $(-7)^{11}$.

\ex Γράψε ως δύναμη: α) $\frac{3}{5}\cdot\frac{3}{5}\cdot\frac{3}{5}$ \quad β) $(-2)(-2)(-2)(-2)$ \quad γ) $0{,}1\cdot0{,}1$ \quad δ) $\left(-\frac{1}{2}\right)\cdot\left(-\frac{1}{2}\right)$.

\ex Υπολόγισε και σύγκρινε: α) $(-2)^4$ και $-2^4$ \quad β) $(-3)^3$ και $-3^3$ \quad γ) $\left(-\frac{1}{2}\right)^2$ και $-\left(\frac{1}{2}\right)^2$.

\ex Υπολόγισε: α) $1^{100}$ \quad β) $0^5$ \quad γ) $(-1)^{100}$ \quad δ) $(-1)^{101}$ \quad ε) $10^3$ \quad στ) $(0{,}1)^2$.

\ex Κάνε τις πράξεις: α) $\left(\frac{1}{2}\right)^2+\left(-\frac{1}{2}\right)^3+\frac{1}{4^2}$ \quad β) $0{,}1^2+\left(\frac{1}{10}\right)^2-0{,}2^2$ \quad γ) $-1^{10}+1^{100}+(-1)^2+(-1)^3$.

\zoni{Το χρησιμοποιώ}{μπάλες, διπλώσεις, εμβαδά}

\ex \textbf{Η μπάλα.} Ο Νίκος αφήνει μια μπάλα από ύψος $2$ μ. και κάθε αναπήδηση φτάνει στο μισό της προηγούμενης. Σε τι ύψος φτάνει μετά την τρίτη και μετά την τέταρτη αναπήδηση;

\ex \textbf{Οι διπλώσεις.} Η Αγγελική διπλώνει ένα φύλλο χαρτί στη μέση πέντε φορές. Τι μέρος του αρχικού εμβαδού έχει το διπλωμένο φύλλο; Γράψε το ως δύναμη και ως κλάσμα.

\ex \textbf{Το τετράγωνο.} Μια αυλή είναι τετράγωνο με πλευρά $1{,}5$ μ. Πόσο είναι το εμβαδόν της; Ο Ηλίας λέει $1{,}5\cdot2=3$. Τι λάθος έκανε;

\ex \textbf{Η παράσταση.} Υπολόγισε $8\cdot\left(-\frac{1}{2}\right)^3+(0{,}5+0{,}3)^2$.

\ex \textbf{Το μοτίβο.} Υπολόγισε τα $\left(\frac{1}{10}\right)^1$, $\left(\frac{1}{10}\right)^2$, $\left(\frac{1}{10}\right)^3$, $\left(\frac{1}{10}\right)^4$ ως δεκαδικούς. Τι παρατηρείς;

\zoni{Το σκέφτομαι}{γιατί βγαίνει έτσι;}

\ex Ο Βασίλης υπολόγισε $-3^2=9$. Εξήγησε το λάθος του και δείξε ότι $(-3)^2=9$.

\ex Πότε είναι θετική μια δύναμη με αρνητική βάση; Εξήγησε με τον κανόνα των προσήμων.

\ex Ποιο είναι μεγαλύτερο, το $\left(\frac{1}{2}\right)^3$ ή το $\left(\frac{1}{2}\right)^2$; Και στη βάση $2$; Τι παρατηρείς για βάσεις μικρότερες και μεγαλύτερες από το $1$;

\ex Αν $\alpha=-\frac{1}{2}$, ποιο είναι μεγαλύτερο, το $\alpha^2$ ή το $\alpha^3$; Εξήγησε χωρίς πράξεις.

\zoni{Για το σπίτι}{πίνακες δυνάμεων}

\ex Φτιάξε πίνακα με τις δυνάμεις $\left(\frac{1}{2}\right)^{\nu}$ για $\nu=1,2,\dots,6$, με κλάσμα και δεκαδικό. Τι παρατηρείς;

\ex Φτιάξε τρεις παραστάσεις με ρητούς και δυνάμεις, που να χρειάζονται και παρενθέσεις, και υπολόγισέ τες.
